\documentclass{article} % For LaTeX2e
\usepackage{iclr2026_conference,times}

\input{math_commands.tex}

\usepackage{hyperref}
\usepackage{url}
\usepackage{amsmath,amssymb,amsfonts,amsthm}
\usepackage{graphicx}
\usepackage{booktabs}
\usepackage{multirow}
\usepackage{microtype}
\usepackage{xcolor}
\usepackage{algorithm}
\usepackage{algpseudocode}
\usepackage{enumitem}

\title{SCALE-Merge: Sign-Consistent and Activation-Aware Linear Equations for Data-Free Model Merging}

\author{Anonymous Authors\\
Anonymous Institution\\
\texttt{anon@example.com}
}

\newcommand{\fix}{\marginpar{FIX}}
\newcommand{\new}{\marginpar{NEW}}
\newcommand{\RESULT}[1]{--}

\newcommand{\dt}{\Delta_t}
\newcommand{\wt}{W_t}
\newcommand{\wo}{W_0}
\newcommand{\ws}{W^\star}

\newtheorem{proposition}{Proposition}

\begin{document}

\maketitle

\begin{abstract}
Model merging combines task-specific fine-tuned checkpoints into a single multi-task model without training. Existing methods split into \emph{data-free} approaches (Task Arithmetic, TIES) that manipulate task vectors $\Delta_t=W_t-W_0$, and \emph{data-dependent} approaches (RegMean, MaTS) that solve a layer-wise least-squares problem using activation covariances. The recent ACTMat method bridged these camps by observing that the activation covariance can be approximated by the gram matrix $\Delta_t^\top\Delta_t$. Starting from the hypothesis that TIES-style interference sources (redundant low-magnitude and sign-conflicting updates) should also degrade this approximation, we propose \textbf{SCALE-Merge} (Sign-Consistent Activation-aware Linear Equations), a fully data-free, closed-form method that first \emph{cleans} task vectors via magnitude-trimming and sign-election and then plugs them into a RegMean-style linear system. SCALE-Merge strictly generalizes ACTMat, Task Arithmetic, and TIES under hyperparameter limits. Across eight image-classification tasks merged into one CLIP ViT-B/32, SCALE-Merge reaches $84.5\%$ average accuracy, outperforming DARE, DELLA, Model Breadcrumbs, Fisher-approx, TIES, Consensus Merging, ACTMat, \emph{and} data-based RegMean by $5.5$--$20$ points; its lead over ACTMat \emph{grows} from parity at $T{=}2$ to $+6.5$~pp at $T{=}8$, and is robust across four orders of magnitude of ridge.
A direct diagnostic against a data-based covariance oracle returns a surprising \emph{negative} result: trimming does \emph{not} make $\Delta_t^\top\Delta_t$ a better approximation of the true activation covariance. Rather, it re-weights the RegMean linear system to concentrate on task-relevant parameter directions. We make this claim quantitative via a \emph{task-routing-matrix} diagnostic showing that SCALE-Merge nearly doubles ACTMat's on/off-target routing concentration. This inductive-bias view of the covariance explains why SCALE-Merge \emph{outperforms} even the data-based RegMean oracle.
\end{abstract}

\section{Introduction}

Fine-tuning large pre-trained models on downstream tasks has become the default recipe for building specialized ML systems \citep{devlin2019bert,dosovitskiy2021image}. This has led to a proliferation of task-specific checkpoints, and a growing interest in recycling them \citep{wortsman2022model,ilharco2023editing}. Model merging \citep{matena2022merging,wortsman2022model,ilharco2023editing} aims to combine $T$ task-specific models $W_1,\ldots,W_T$ (all fine-tuned from a common pre-trained model $W_0$) into a single multi-task model $\ws$ \emph{without access to the training data and without extra training}. Compared to joint multi-task training, merging requires only the final weights; compared to output-space ensembling, it produces a model that is $T$-times cheaper to evaluate.

A useful abstraction is the \emph{task vector} \citep{ilharco2023editing}: $\Delta_t = W_t - W_0$. Early data-free methods---Task Arithmetic \citep{ilharco2023editing} and TIES-Merging \citep{yadav2023tiesmerging}---construct $\ws$ as $W_0 + \alpha \sum_t \tilde\Delta_t$, where $\tilde\Delta_t$ is either the raw task vector or a sparsified/sign-corrected version. In parallel, data-dependent approaches such as RegMean \citep{jin2023dataless} and MaTS \citep{tam2024mats} solve a per-layer least-squares problem,
\begin{equation}
\ws = \left(\sum_t C_t\right)^{\dagger} \sum_t C_t\, W_t,
\label{eq:regmean}
\end{equation}
where $C_t \in \mathbb{R}^{d_{\text{in}} \times d_{\text{in}}}$ is the activation covariance on task $t$'s training data. Solving~\eqref{eq:regmean} with accurate $C_t$'s matches activations and typically yields the best multi-task accuracy, but requires access to per-task input data at merging time---a significant practical limitation when merging checkpoints shared on the open Internet.

The recent ACTMat method \citep{hameed2025actmat} closes this gap by showing that the second-moment matrix of the inputs can be approximated directly from the task vector: $C_t \approx \Delta_t^\top \Delta_t$. This yields a \emph{fully data-free} closed form that enjoys much of the performance of data-based RegMean. Our starting hypothesis was that ACTMat's bound on $\angle(\Delta_t^\top \Delta_t, C_t)$---which decomposes into cross-term, correlation, and drift terms that all deteriorate when $\Delta_t$ contains small-magnitude or sign-flipped entries---should be tightened by the same TIES-style preprocessing \citep{yadav2023tiesmerging} that removes these interference sources. We therefore built SCALE-Merge, and found that it substantially outperforms ACTMat. Surprisingly, however, a direct empirical check (\S\ref{sec:analysis}) shows that trim-and-elect preprocessing does \emph{not} improve the angular approximation of $C_t$---instead it degrades it. This null finding reframes the paper: SCALE-Merge wins not because it better approximates the data-based covariance, but because the \emph{cleaned task-vector gram matrix} is a better effective reweighting for the merge linear system. We make this reframing concrete below (\S\ref{sec:analysis}) by comparing against a data-based RegMean oracle that SCALE-Merge nonetheless outperforms.

\textbf{Our contribution.} We propose \textbf{SCALE-Merge} (\textbf{S}ign-\textbf{C}onsistent and \textbf{A}ctivation-aware \textbf{L}inear \textbf{E}quations for merging), which unifies these two perspectives. Given task-specific fine-tuned weights $W_1,\ldots,W_T$ and the pre-trained weights $W_0$, SCALE-Merge (i) trims each task vector by keeping only the top-$k\%$ of entries by magnitude, (ii) elects a consensus sign per parameter and zeroes out entries of $\Delta_t$ that disagree with the consensus, and (iii) uses the cleaned task vectors $\tilde\Delta_t$ as the covariance $\hat C_t = \tilde\Delta_t^\top \tilde\Delta_t$ in a RegMean-style per-layer linear system. The resulting merging procedure is fully data-free, closed-form, and strictly generalizes ACTMat (which corresponds to SCALE with $k=100\%$ and no sign election).

\begin{figure}[t]
\centering
\includegraphics[width=\linewidth]{../results/method_diagram.png}
\caption{\textbf{SCALE-Merge pipeline} (per linear weight). Each task's task vector $\Delta_t = W_t - W_0$ is trimmed by magnitude (TIES step 1) and sign-aligned with a global consensus (TIES step 2). The \emph{cleaned} task vectors $\tilde\Delta_t$ serve as a data-free estimate of the input covariance $\hat C_t \approx \mathbb{E}_{z\sim\mathcal{D}_t}[zz^\top]$, which is plugged into the RegMean/MaTS closed-form linear system to obtain the merged weight $W^\star$.}
\label{fig:diagram}
\end{figure}

We make the following contributions:
\begin{itemize}
    \item We introduce \textbf{SCALE-Merge}, a simple, closed-form, data-free merging rule that unifies TIES-style interference resolution with activation-matching covariance merging, strictly generalizing ACTMat and recovering Simple Averaging, Task Arithmetic, and TIES as special cases.
    \item On 8 image-classification tasks merged into a CLIP ViT-B/32, SCALE-Merge outperforms Task Arithmetic, TIES, DARE, DELLA, Model Breadcrumbs, Fisher-approx, Consensus Merging, ACTMat, \emph{and} data-based RegMean (which uses true activation covariances) by a meaningful margin in average accuracy. SCALE-Merge's lead over ACTMat \emph{grows} with the number of merged tasks---from parity at $T{=}2$ to $+6.5$~pp at $T{=}8$---and the optimal trim rate $k^\star$ decreases monotonically with $T$.
    \item We present a \emph{negative empirical result} that reshapes the theoretical picture: TIES-style trimming does \emph{not} reduce the angular distance between $\Delta_t^\top\Delta_t$ and the data-based $C_t$. The cleaned gram matrix is a worse estimator of $C_t$, yet a better effective reweighting for the RegMean linear system---explaining why SCALE-Merge beats even data-based RegMean.
    \item We introduce a \emph{task-routing-matrix} diagnostic that quantifies what effective reweighting trim-and-elect provides: SCALE-Merge concentrates per-task routing mass on task-relevant parameter columns, raising the on/off-target ratio from $1.24$ (ACTMat) to $1.42$ at $k{=}0.3$ and $2.24$ at $k{=}0.05$. An ablation further shows that trimming dominates sign-election on this benchmark because top-$k$ task-vector supports are already near-disjoint at $k\leq 0.3$ (agreement rate $\to 1$).
\end{itemize}

\section{Related Work}
\label{sec:related}

\paragraph{Model merging and weight averaging.} Early work on weight averaging \citep{mcmahan2017communication,stich2018local,izmailov2018averaging,neyshabur2020being,frankle2020linear} showed that simple averaging of weights from independent training runs can improve generalization when the models share initialization. \citet{wortsman2022model} showed the same holds across fine-tuned checkpoints of a common pre-trained model, introducing Model Soups. \citet{matena2022merging} proposed Fisher-weighted averaging to account for parameter importance, and \citet{daheim2024model} extended this to gradient-matching under uncertainty. \citet{ilharco2023editing} introduced Task Arithmetic as a way to edit pre-trained models in weight space by adding and subtracting \emph{task vectors}. Later work has explored alternative parameterizations \citep{ortiz2023task}, weight disentanglement \citep{ortiz2024weight,jin2024fine}, permutation alignment \citep{singh2020neuron,stoica2024zipit}, and safety-preserving merging \citep{bhardwaj2024language}. Practical toolkits such as MergeKit \citep{goddard2024arcee} and surveys \citep{yang2024model} reflect the rapid maturation of the field.

\paragraph{Sparsification-based interference resolution.} \citet{yadav2023tiesmerging} introduced TIES-Merging, which trims redundant parameters and resolves sign conflicts before averaging. \citet{yu2024language} study random drop-and-rescale (DARE), \citet{deep2024della} propose magnitude-aware Bernoulli sampling (DELLA), and \citet{davari2024model} introduce Model Breadcrumbs (drop smallest and largest entries of $\Delta_t$). These four methods all ask the same first-order question---\emph{which entries of $\Delta_t$ should survive aggregation?}---and our experiments (\S\ref{sec:exp}) show this family plateaus below $70\%$ on the 8-task ViT-B/32 benchmark. \citet{adamerging2024,adamerging2024adp,adarank2025,akiba2024evolutionary,huang2024emr} combine these ideas with adaptive weighting, rank pruning, or evolutionary search, but still without activation-covariance information.

\paragraph{Subspace and support localization.} A second, largely orthogonal line of work asks not which entries but \emph{which directions} of $\Delta_t$ carry task information. \citet{panigrahi2023task} localize task competence in a small set of neurons; \citet{wang2024localizing} formalize this with TALL-masks and Consensus Merging, which retains only entries where $\geq q$ of the $T$ tasks agree that an entry is important (we include it as a baseline in \S\ref{sec:exp}). \citet{marczak2025tsv} introduce Task Singular Vectors to reduce inter-task subspace overlap, and \citet{tang2023concrete} learn a merging subspace end-to-end. These methods are conceptually close to our \emph{task-routing-matrix} diagnostic (\S\ref{sec:analysis}); SCALE-Merge differs in that its routing matrix emerges \emph{implicitly} from the cleaned gram matrix, with no explicit mask or auxiliary subspace, yet achieves on/off-target concentration nearly $2\times$ higher than ACTMat, echoing evidence that a small parameter subspace in pre-trained models carries disproportionate functional mass \citep{geva2021transformer,molchanov2017variational}.

\paragraph{Activation-matching methods.} A separate line of work formulates merging as layer-wise activation matching: $\ws = \arg\min_W \sum_t \mathbb{E}_{z\sim\mathcal{D}_t}\|Wz - W_t z\|^2_2$, with closed-form solution~\eqref{eq:regmean} \citep{jin2023dataless}. \citet{tam2024mats} unify simple averaging, Fisher merging, and RegMean under a \emph{task parameter subspace} perspective and use conjugate gradient to scale the framework. \citet{regmean2025} improve RegMean with better regularization. \citet{aceemrging2026} adaptively combine multiple covariance estimators. These methods rely on access to per-task input data at merge time, making them impractical in the growing ``model hub'' setting where checkpoints are downloaded without their training data.

\paragraph{Data-free activation matching.} \citet{hameed2025actmat} observe that the input covariance $C_t = \mathbb{E}_{z\sim\mathcal{D}_t}[zz^\top]$ is approximately proportional to the gram matrix of the task vector, $\Delta_t^\top \Delta_t$, when the layer is fine-tuned from $W_0$ by gradient descent. This yields a fully data-free merging rule, which they call ACTMat. Our work builds directly on ACTMat: we observe that its data-free covariance estimate is susceptible to the two types of interference that TIES identifies, and propose a simple preprocessing step that sharpens each ACTMat angular error term. Concurrent work by \citet{bowen2024weighted} also explores data-free covariance weighting, but does not incorporate sign-based interference resolution.

\paragraph{Other directions.} Merging has also been studied for diffusion models \citep{nair2024diffmerge}, LoRA adapters \citep{zhang2024loremerging,huang2024lorahub}, continual learning \citep{continualmerge2025,nullspace2025}, multimodal models \citep{sung2023empirical}, and LLM safety-alignment \citep{bhardwaj2024language}. \citet{gueta2023knowledge} study the geometry of fine-tuned weights. \citet{twinmerging2024} combine shared and task-specific parts. We focus on the data-free multi-task merging setting common to TIES, ACTMat, and MaTS.

\section{Preliminaries}

\paragraph{Notation.} Let $W_0$ denote the pre-trained model's weights and $W_1,\ldots,W_T$ be weights of $T$ task-specific fine-tuned models obtained by further training $W_0$ on tasks $1,\ldots,T$ respectively. The task vector for task $t$ is $\Delta_t = W_t - W_0$. For a linear layer with input $z \in \mathbb{R}^{d_\text{in}}$, the activation second-moment matrix is $C_t = \mathbb{E}_{z\sim\mathcal{D}_t}[z z^\top] \in \mathbb{R}^{d_\text{in}\times d_\text{in}}$.

\paragraph{RegMean / MaTS.} The activation-matching objective is $\ws = \arg\min_W \sum_t \|(W-W_t)z\|^2_{z\sim \mathcal{D}_t}$, whose closed-form solution is~\eqref{eq:regmean}. Simple averaging, Fisher merging, and RegMean differ only in the choice of $C_t$ \citep{tam2024mats}.

\paragraph{ACTMat.} \citet{hameed2025actmat} prove that, under mild conditions (full-batch gradient descent, fixed learning rate, bounded drift), the angular distance between $\Delta_t^\top\Delta_t$ and $C_t^{(K)}$ is bounded by the sum of three terms: a \emph{cross-term error} $\epsilon^{(\text{cross})} = \angle(G^\top G, S)$; a \emph{correlation error} $\epsilon^{(\text{corr})} = \angle(S, \tilde S)$; and a \emph{drift error} $\epsilon^{(\text{drift})} = \angle(\tilde S, C_t^{(K)})$, where $G,S,\tilde S$ are defined in \citet[Theorem 3.1]{hameed2025actmat}. All three terms shrink as the gradients and activations become more concentrated on a small set of informative directions.

\paragraph{TIES-Merging.} \citet{yadav2023tiesmerging} empirically document two sources of interference: (i) \emph{redundant} low-magnitude entries in $\Delta_t$ that do not carry useful task information but add noise when summed across tasks, and (ii) \emph{sign-conflicting} entries where tasks disagree about the sign of a parameter's update. TIES addresses both: it trims each $\Delta_t$ to keep only its top-$k\%$ of entries by magnitude, then elects a global sign per parameter and averages only over tasks whose sign agrees with the elected sign.

\section{Method: SCALE-Merge}
\label{sec:method}

The key observation that motivates SCALE-Merge is:
\begin{quote}\emph{The two interference patterns TIES identifies in $\Delta_t$ are the exact sources of error in ACTMat's data-free covariance estimate.}\end{quote}

Concretely: (a) low-magnitude entries of $\Delta_t$ contribute mostly noise to the off-diagonal of $\Delta_t^\top \Delta_t$, inflating $\epsilon^{(\text{cross})}$ and $\epsilon^{(\text{corr})}$; (b) sign-flipped entries across tasks inflate the \emph{cross-task} covariance estimate $\sum_t \Delta_t^\top \Delta_t$ without corresponding shift in the cross-task sum $\sum_t C_t$. We therefore propose to apply TIES-style preprocessing to $\Delta_t$ \emph{before} forming the covariance estimate.

\paragraph{Algorithm.} For each 2-D weight $W\in\mathbb{R}^{d_\text{out}\times d_\text{in}}$ (higher-rank weights reshaped accordingly), SCALE-Merge applies four steps: \textbf{(i) Trim} each task vector $\Delta_t$ to its top-$k\%$ entries by magnitude ($\hat\Delta_t$); \textbf{(ii) Elect sign} $\gamma=\mathrm{sgn}(\sum_t \hat\Delta_t)$ and zero entries whose sign disagrees with $\gamma$ (yielding $\tilde\Delta_t$); \textbf{(iii) Form covariance} $\hat C_t = \tilde\Delta_t^\top\tilde\Delta_t$; \textbf{(iv) Solve} $\ws = (\sum_t\hat C_t + \rho\bar\lambda I)^{-1}\sum_t W_t\hat C_t$ with $\bar\lambda = \tfrac{1}{d_\text{in}}\mathrm{tr}(\sum_t\hat C_t)$. One-dimensional parameters (LayerNorm scales, positional embeddings) fall back to simple averaging. Pseudocode in Algorithm~\ref{alg:scale}.

\begin{algorithm}[t]
\caption{SCALE-Merge (per linear weight)}
\label{alg:scale}
\begin{algorithmic}[1]
\Require Pretrained weight $W_0 \in \mathbb{R}^{d_\text{out}\times d_\text{in}}$; task-specific weights $\{W_t\}_{t=1}^T$; keep-fraction $k$; ridge $\rho$.
\For{$t = 1,\ldots,T$}
  \State $\Delta_t \gets W_t - W_0$ \Comment{task vector}
  \State $m_t \gets$ mask keeping top-$k\%$ of $|\Delta_t|$ entries, else $0$ \Comment{trim}
  \State $\hat\Delta_t \gets \Delta_t \odot m_t$
\EndFor
\State $\gamma \gets \mathrm{sgn}\!\left(\sum_t \hat\Delta_t\right)$ \Comment{elect sign}
\For{$t = 1,\ldots,T$}
  \State $\tilde\Delta_t \gets \hat\Delta_t \odot \mathbb{1}[\mathrm{sgn}(\hat\Delta_t) = \gamma]$ \Comment{zero disagreeing entries}
  \State $\hat C_t \gets \tilde\Delta_t^\top \tilde\Delta_t \in \mathbb{R}^{d_\text{in}\times d_\text{in}}$ \Comment{data-free covariance}
\EndFor
\State $C \gets \sum_t \hat C_t + \rho\cdot \bar\lambda\cdot I$ where $\bar\lambda = \tfrac{1}{d_\text{in}}\mathrm{tr}(\sum_t \hat C_t)$
\State $W^\star \gets \left(\sum_t W_t\, \hat C_t\right) C^{-1}$ \Comment{closed-form solve}
\State \Return $W^\star$
\end{algorithmic}
\end{algorithm}

\paragraph{Special cases.} SCALE-Merge generalizes several prior methods by varying $(k, \text{sign}, \rho)$:
\begin{itemize}[leftmargin=1.5em]
\item $k=1.0$, no sign election, $\rho\to0$: recovers \textbf{ACTMat} \citep{hameed2025actmat}.
\item $k=1.0$, no sign election, $\hat C_t = I$: recovers \textbf{simple averaging} (because $\sum_t W_t$ is returned up to normalization).
\item $k \in (0,1)$, sign election, $\hat C_t = I$ replacing the covariance by a uniform weight: recovers a covariance-free variant close to \textbf{TIES} (see~\S\ref{sec:exp} for the exact comparison).
\end{itemize}
Thus SCALE-Merge interpolates between Task Arithmetic / TIES (which do not exploit second-order structure) and ACTMat / RegMean (which do, but trust noisy data).

\paragraph{Initial motivation (partly wrong, but still productive).} Our original motivation was to tighten ACTMat's angular-distance bound on $\angle(\Delta_t^\top\Delta_t, C_t)$: the three error terms in \citet[Theorem 3.1]{hameed2025actmat} depend on a double sum of the form $\sum_{k,k'}\mathbb{E}[z^{(k)} z^{\prime (k')\top} g^{(k)\top} g^{\prime(k')}]$, and zero-magnitude or sign-conflicted entries of $\Delta_t$ correspond to gradient-activation products that are either negligibly small or noise-dominated. Removing them therefore \emph{seemed} likely to reduce all three error terms. In \S\ref{sec:analysis} we test this hypothesis directly against a data-based oracle and find it is \emph{not} supported: trimming does not tighten the ACTMat approximation. Nevertheless, SCALE-Merge substantially outperforms both ACTMat and data-based RegMean (\S\ref{sec:exp}). We revise our theoretical picture accordingly.

\paragraph{Revised interpretation: SCALE as effective reweighting.} The RegMean/MaTS solve is $\ws = (\sum_t \hat C_t)^{-1} \sum_t \hat C_t W_t$. What enters this system is not any individual $\hat C_t$'s \emph{approximation quality} of $C_t$, but its \emph{relative weight} across parameter directions. In that sense the $\hat C_t$'s in ACTMat/SCALE are better thought of as \emph{task-relevance weightings} than as covariance estimates. Trimming concentrates the weight onto parameter directions where task $t$ moved far from the pre-trained weights---exactly the directions where borrowing from $W_t$ is most beneficial and least interfering with other tasks. This is the same spirit as TIES's trim step, but extended from the first-order sum $\sum_t\tilde\Delta_t$ to the second-order weighting $\sum_t\tilde\Delta_t^\top\tilde\Delta_t$. Sign-election plays a complementary role for the diagonal entries (writing out $(i,j)$ with $i\neq j$: $\Delta_1[i]\Delta_1[j] + \Delta_2[i]\Delta_2[j]$ can cancel destructively when signs flip), though on our benchmark trimming already renders top-$k$ supports near-disjoint (\S\ref{sec:analysis}).

\paragraph{Computational cost.} SCALE-Merge performs one matrix solve of size $d_\text{in}\times d_\text{in}$ per layer and $O(T\cdot d_\text{out}\cdot d_\text{in}^2)$ matmuls for the gram matrices. For ViT-B/32, $d_\text{in} \le 3072$, so the full merge runs in seconds on a single GPU---comparable to ACTMat and much cheaper than any RegMean variant that requires forward passes.

\section{Experiments}
\label{sec:exp}

\subsection{Setup}

\paragraph{Base model.} We use CLIP ViT-B/32 (OpenAI pre-trained) \citep{radford2021learning} via the \texttt{open\_clip} library \citep{openclip2023}. Following the standard protocol of \citet{ilharco2023editing,yadav2023tiesmerging}, we fine-tune only the visual encoder and use frozen CLIP text embeddings as the classifier head (one head per task). Only visual-encoder parameters are merged.

\paragraph{Tasks.} We fine-tune CLIP ViT-B/32 on \textbf{eight image classification tasks}: MNIST \citep{lecun1998mnist}, CIFAR-10 \citep{krizhevsky2009cifar}, CIFAR-100, SVHN \citep{netzer2011svhn}, FashionMNIST \citep{xiao2017fashion}, EuroSAT \citep{helber2019eurosat}, GTSRB \citep{stallkamp2011gtsrb}, and DTD \citep{cimpoi2014dtd}. We use AdamW~\citep{loshchilov2019decoupled} in PyTorch~\citep{paszke2019pytorch}, cosine LR schedule with $10\%$ warm-up, peak LR $10^{-5}$, weight decay $0.1$, batch size $128$, for up to $5{-}10$ epochs per task.

\paragraph{Baselines.} We compare against:
\begin{itemize}[leftmargin=1.5em]
\item \textbf{Pre-trained zero-shot}: evaluate $W_0$ with CLIP text classifiers.
\item \textbf{Individual fine-tuned} (upper bound): evaluate each $W_t$ on its own task.
\item \textbf{Simple Averaging}: $\ws = \tfrac{1}{T}\sum_t W_t$.
\item \textbf{Task Arithmetic} \citep{ilharco2023editing}: $\ws = W_0 + \alpha \sum_t \Delta_t$, $\alpha$ swept.
\item \textbf{TIES-Merging} \citep{yadav2023tiesmerging}: trim-elect-merge with $k=0.2$, $\alpha$ swept.
\item \textbf{DARE} \citep{yu2024language}: random drop-and-rescale (drop rate $p{=}0.9$) applied to task vectors, then summed as Task Arithmetic; $\alpha$ swept.
\item \textbf{DARE+TIES}: DARE sparsification followed by TIES sign-election and disjoint merge.
\item \textbf{DELLA} \citep{deep2024della}: magnitude-ranked Bernoulli sampling where low-magnitude entries of $\Delta_t$ are kept with probability $p_\text{low}$ and high-magnitude entries with $p_\text{high}$ (linear interpolation across rank), survivors are rescaled by $1/p$, and the resulting task vectors are aggregated via TIES elect-sign. We sweep $p_\text{low}\in\{0.05,0.1,0.2\}$, $p_\text{high}\in\{0.8,0.9\}$, $\alpha\in\{0.2,0.3,0.5\}$ and report the best configuration.
\item \textbf{Model Breadcrumbs} \citep{davari2024model}: deterministic sparsification that zeroes both the smallest-magnitude and largest-magnitude entries of each $\Delta_t$, keeping only the middle band. We sweep $(\text{drop-small}\in\{0.85,0.9,0.92\},\ \text{drop-large}\in\{0.005,0.01,0.02\},\ \alpha\in\{0.2,0.3,0.5\})$ and report the best configuration.
\item \textbf{Fisher-approx} \citep{matena2022merging}: data-free Fisher-weighted averaging, with per-parameter Fisher diagonals approximated by $\Delta_t^2$ (Pascanu \& Bengio 2014 approximation in the data-free regime).
\item \textbf{ACTMat} \citep{hameed2025actmat}: data-free covariance merging with $C_t = \Delta_t^\top\Delta_t$.
\item \textbf{Consensus Merging (TALL-masks)} \citep{wang2024localizing}: a data-free baseline that first builds a per-task importance mask $m_t = \mathbb{1}[|\Delta_t|\in \text{top-}k\%]$, intersects the masks with a minimum-agreement threshold $q$ to obtain a consensus mask $M = \mathbb{1}[\sum_t m_t \geq q]$, and merges as $\ws = W_0 + \alpha\, M \odot \sum_t \Delta_t$; $(k,q,\alpha)$ swept.
\item \textbf{RegMean (oracle, data-based)} \citep{jin2023dataless}: we additionally run a data-based RegMean baseline in which the per-task input-activation covariance at every linear layer is computed by running $1{,}024$ task-$t$ training images through the fine-tuned $W_t$. This is \emph{not} data-free; we include it as an upper-bound reference for covariance-based methods on our benchmark.
\end{itemize}

\paragraph{Metric.} Per-task top-1 accuracy on the official test split, and the average across tasks. For EuroSAT we use an 80/20 random split of the training data since the dataset ships without a held-out test set.

\subsection{Main results}

\begin{table}[t]
\centering
\caption{Top-1 test accuracy per task (\%) and average across the 8 tasks, when merging 8 task-specific CLIP ViT-B/32 checkpoints into a single visual encoder. \textbf{Individual} reports each $W_t$ evaluated on its own task (an upper bound; it is not itself a merged model). \textbf{Pretrained} is zero-shot CLIP. Best merged result per column in bold.}
\label{tab:main}
\footnotesize
\setlength{\tabcolsep}{4pt}
\resizebox{\linewidth}{!}{
\begin{tabular}{lccccccccc}
\toprule
Method & MNIST & CIFAR10 & CIFAR100 & SVHN & FashionM & EuroSAT & GTSRB & DTD & \textbf{Avg} \\
\midrule
Pretrained (zero-shot) & 30.3 & 88.8 & 61.7 & 9.3 & 61.5 & 33.0 & 26.5 & 43.2 & 44.3 \\
Individual (UB)        & 99.6 & 97.7 & 89.0 & 97.3 & 95.0 & 98.7 & 98.7 & 75.0 & 93.9 \\
\midrule
Simple Averaging       & 69.1 & 94.2 & 72.0 & 33.6 & 75.0 & 57.1 & 44.0 & 48.3 & 61.7 \\
Task Arithmetic \citep{ilharco2023editing}       & 82.6 & 94.0 & 70.0 & 44.6 & 77.1 & 58.4 & 50.1 & 46.9 & 65.5 \\
DARE \citep{yu2024language} & 82.8 & 94.0 & 69.9 & 44.5 & 77.4 & 57.4 & 50.1 & 46.8 & 65.3 \\
TIES-Merging \citep{yadav2023tiesmerging}          & 77.8 & 95.3 & 74.9 & 42.2 & 76.1 & 53.5 & 47.5 & 44.9 & 64.0 \\
DARE + TIES            & 85.6 & 93.6 & 68.7 & 48.6 & 78.1 & 52.9 & 51.9 & 43.6 & 65.4 \\
DELLA \citep{deep2024della} & 84.1 & 95.0 & 73.0 & 48.1 & 77.1 & 52.3 & 49.9 & 42.9 & 65.3 \\
Model Breadcrumbs \citep{davari2024model} & 87.4 & 94.0 & 69.2 & 49.0 & 78.1 & 57.8 & 53.6 & 47.7 & 67.1 \\
Consensus Merging (TALL-masks) \citep{wang2024localizing} & 87.4 & 93.5 & 67.1 & 50.1 & 77.9 & 55.9 & 52.8 & 45.6 & 66.3 \\
Fisher-approx \citep{matena2022merging}         & 72.1 & 94.5 & 72.5 & 36.3 & 75.4 & 57.0 & 44.8 & 47.2 & 62.5 \\
ACTMat \citep{hameed2025actmat}                & 99.0 & 92.7 & 64.8 & 84.3 & 89.6 & 72.5 & 85.5 & 35.5 & 78.0 \\
\midrule
RegMean (\emph{data-based}) \citep{jin2023dataless}  & 96.7 & 96.2 & 77.0 & 63.2 & 87.4 & 92.6 & 57.4 & 61.5 & 79.0 \\
\midrule
\textbf{SCALE-Merge (ours, data-free)}  & \textbf{99.0} & \textbf{96.1} & \textbf{76.1} & 83.2 & 89.5 & \textbf{91.4} & \textbf{86.4} & \textbf{54.6} & \textbf{84.5} \\
\bottomrule
\end{tabular}}
\end{table}

Table~\ref{tab:main} and Figure~\ref{fig:avg} report merged accuracy on all eight tasks. SCALE-Merge reaches $\mathbf{84.5\%}$ average accuracy, improving over the best data-free baseline (ACTMat, $78.0\%$) by $6.5$ points and over Task Arithmetic ($65.5\%$) by $19.0$ points. The merged model closes $81\%$ of the zero-shot~$\to$~individual-fine-tuned gap ($93.9 - 44.3 = 49.6$ points, of which SCALE-Merge recovers $40.2$). Per-task, SCALE-Merge is the \emph{top} merger (including data-based RegMean) on 6 of 8 tasks (Figure~\ref{fig:per_task}); on SVHN and FashionMNIST it is within $1$ point of ACTMat. The improvement is particularly large on tasks where the pretrained CLIP accuracy is poor (e.g.\ $+23\%$ on EuroSAT, $+19\%$ on DTD over ACTMat), indicating that SCALE-Merge effectively injects task-specific knowledge without destructive interference.

\paragraph{Comparison against sparsification baselines (DELLA, Breadcrumbs, Consensus).} DELLA \citep{deep2024della} replaces DARE's uniform Bernoulli drop with a \emph{magnitude-ranked} Bernoulli drop (higher keep-probability for high-magnitude $\Delta_t$ entries) and aggregates via TIES elect-sign; after sweeping $p_\text{low}{\in}\{0.05,0.1,0.2\}$, $p_\text{high}{\in}\{0.8,0.9\}$, $\alpha{\in}\{0.2,0.3,0.5\}$ (18 configs) the best reaches $\mathbf{65.3\%}$. Breadcrumbs \citep{davari2024model} keeps only a middle magnitude band of $\Delta_t$; best over drop-small${\in}\{0.85,0.9,0.92\}$, drop-large${\in}\{0.005,0.01,0.02\}$, $\alpha{\in}\{0.2,0.3,0.5\}$ reaches $\mathbf{67.1\%}$. Consensus Merging / TALL-masks \citep{wang2024localizing} keeps only entries on which at least $q$ of $T$ tasks agree; best over $(k{\in}\{0.1,0.2,0.3\}, q{\in}\{2,3,4\}, \alpha{\in}\{0.2,0.3,0.5\})$ reaches $\mathbf{66.3\%}$. All three first-order sparsification methods cluster near TIES/DARE ($64$--$67\%$), well below ACTMat ($78.0$) and SCALE-Merge ($84.5$). Consensus in particular is consistent with our sign-agreement measurement (\S\ref{sec:exp}): at $k{\leq}0.3$ the top-$k$ supports across tasks are near-disjoint, so the intersection is mostly empty and discards most task-specific mass. None of these methods exploit activation information; the gap between them and SCALE-Merge localizes the gain to the activation-aware quadratic form $\tilde\Delta_t^\top\tilde\Delta_t$, not to sparsification or sign resolution per se.

\paragraph{Comparison against data-based RegMean.} Strikingly, SCALE-Merge also beats the data-based RegMean \emph{oracle} (which uses true activation covariances $C_t$ estimated from 1{,}024 task-$t$ training images through $W_t$) by $5.5$ points of average accuracy ($84.5$ vs.\ $79.0$). The gap is largest on tasks where the task-vector magnitudes are large (SVHN: $+20.0$, GTSRB: $+29.0$) and negative on DTD ($-6.9$; SCALE's trim excludes informative DTD-specific entries that data-based RegMean's true covariance still picks up). A ridge sweep for RegMean over $\rho\in\{10^{-6},10^{-5},10^{-4},10^{-3},10^{-2}\}$ produces a narrow $78.3$--$79.1\%$ range, so this gap is robust and not a tuning artifact. The fact that the data-based oracle \emph{loses} to the trimmed task-vector gram matrix is what motivates the diagnostic in \S\ref{sec:analysis}.

\begin{figure}[t]
\centering
\includegraphics[width=0.7\linewidth]{../results/avg_accuracy.png}
\caption{Average test accuracy across 8 tasks (CLIP ViT-B/32 backbone). SCALE-Merge outperforms all other data-free merging baselines.}
\label{fig:avg}
\end{figure}

\begin{figure}[t]
\centering
\includegraphics[width=\linewidth]{../results/per_task.png}
\caption{Per-task test accuracy for each merging method.}
\label{fig:per_task}
\end{figure}

\subsection{Diagnostic: does trimming improve the ACTMat approximation?}
\label{sec:analysis}

Our original motivation for SCALE-Merge (\S\ref{sec:method}) was that TIES-style trimming and sign-election should tighten ACTMat's bound on $\angle(\Delta_t^\top\Delta_t, C_t)$. We test this hypothesis \emph{directly} by measuring the angular distance to a data-based $C_t$.

\paragraph{Setup.} For every task $t\in\{1,\ldots,8\}$ and for a representative set of $11$ linear layers spanning the depth of the CLIP ViT-B/32 visual encoder, we compute the true input-activation covariance $C_t = \frac{1}{N}\sum_{n=1}^N z_n z_n^\top$ by passing $N{=}256$ task-$t$ training images through $W_t$ and logging per-layer inputs via forward hooks. We then compare the angular distance
$\angle(A, B) = \arccos\bigl(\langle A, B\rangle_F / (\|A\|_F\|B\|_F)\bigr)$
between $C_t$ (resp.\ $\sum_t C_t$) and the data-free gram estimate for a range of keep-fractions $k\in\{0.05, 0.1, 0.2, 0.3, 0.5, 1.0\}$, with and without sign-election. We also report the top-$32$ eigenspace overlap (the average squared projection of the top-$32$ eigenvectors of one matrix onto the top-$32$ eigenspace of the other; 1~= perfect overlap), which is arguably the quantity that matters for the RegMean linear solve since low-eigenvalue directions contribute little.

\paragraph{Result (negative).} Figure~\ref{fig:angular} summarizes the measurements, averaged across layers and tasks. Trimming the task vector \emph{does not} reduce the angular error---it \emph{increases} it slightly, both per-task ($74.1^\circ \to 76.3^\circ$ as $k$ drops from $1.0$ to $0.05$) and at the summed level ($73.6^\circ \to 75.1^\circ$). The top-$32$ subspace overlap \emph{also} drops with trimming (from $0.33$ at $k{=}1$ to $0.29$ at $k{=}0.05$). Sign-election barely moves the angle ($\leq 0.25^\circ$). Our hypothesis that SCALE tightens the ACTMat angular bound is thus empirically refuted on this benchmark.

\begin{figure}[t]
\centering
\includegraphics[width=0.95\linewidth]{../results/angular_distance.png}
\caption{Angular distance (left) and top-$32$ eigenspace overlap (right) between the data-free gram estimate $\Delta_t^\top\Delta_t$ and the data-based covariance $C_t$, averaged over $11$ representative linear layers and 8 tasks. Both quantities \emph{degrade} as the trim keep-fraction $k$ decreases, showing that trim-and-elect does not improve the underlying covariance approximation. Sign-election barely moves these curves (not shown; $\leq 0.25^\circ$). Yet trim-and-elect \emph{does} substantially improve merged accuracy (Table~\ref{tab:ablation}). We conclude that SCALE-Merge's gains come from a better \emph{effective reweighting} of the RegMean linear system rather than a better approximation of $C_t$---consistent with SCALE-Merge beating even the data-based RegMean (Table~\ref{tab:main}).}
\label{fig:angular}
\end{figure}

\paragraph{So why does SCALE-Merge work?} We interpret the gain as an \emph{effective reweighting} of the RegMean linear system rather than better covariance estimation---corroborated by three observations: (a) SCALE beats data-based RegMean (Table~\ref{tab:main}), (b) the gain over ACTMat grows with the number of tasks (Figure~\ref{fig:taskcount}), consistent with reweighting mattering most under heavy cross-task competition, and (c) trim-only ($+6.5$pp) exceeds sign-only ($+4.0$pp), matching the larger weighting shift that trimming produces. Proposition~\ref{prop:rayleigh} below makes this reweighting concrete at the column level.

\paragraph{A routing-matrix diagnostic.} To make the ``effective reweighting'' view \emph{quantitative} we study SCALE-Merge's per-layer \emph{task-routing matrix}
\begin{equation}
P_t(k) \;=\; \hat C_t(k)\,\Bigl(\sum\nolimits_s \hat C_s(k) + \rho\bar\lambda I\Bigr)^{-1}\,\in\,\mathbb{R}^{d_\text{in}\times d_\text{in}}.
\label{eq:routing}
\end{equation}
Written in closed form, the merged weight per layer is $\ws = \sum_t W_t\, P_t(k)$: $P_t$ is the linear operator that decides, column-by-column, how much of task $t$'s fine-tuned weight $W_t$ enters the merged solution. Under this decomposition, a good merge should have each $P_t$ concentrate its mass on \emph{task-$t$ relevant} columns (where $W_t$ moved away from $W_0$), and little mass on columns dominated by other tasks.

\begin{proposition}[Column-wise Fisher view of SCALE-Merge]
\label{prop:rayleigh}
Let $\hat C_t = \tilde\Delta_t^\top\tilde\Delta_t$ and $M = \sum_s \hat C_s + \rho\bar\lambda I$. For any input column $j\in[d_\text{in}]$, SCALE-Merge satisfies $\ws[:,j] = \sum_t W_t\,P_t[:,j]$ with $P_t[:,j] = \hat C_t\, M^{-1} e_j$. When $\{\hat C_t\}$ and $M$ are simultaneously diagonal in the coordinate basis---exact if column supports of $\{\tilde\Delta_t\}$ are pairwise disjoint, and approximate otherwise---the diagonal routing weight reduces to
\begin{equation}
P_t[j,j] \;=\; \frac{\|\tilde\Delta_t[:,j]\|_2^2}{\sum_s \|\tilde\Delta_s[:,j]\|_2^2 \;+\; \rho\bar\lambda}.
\label{eq:fisher-col}
\end{equation}
\end{proposition}

\noindent\emph{Proof sketch.} Using $\hat C_t[j,j]=\|\tilde\Delta_t[:,j]\|_2^2$ and $M^{-1}e_j = e_j / M[j,j]$ in the diagonal limit gives~\eqref{eq:fisher-col}; off-diagonals of $\hat C_t$ (inner products between distinct columns of $\tilde\Delta_t$) vanish when those columns are disjointly supported, which the sign-agreement measurement (\S\ref{sec:exp}, $0.65 \to 0.95$ as $k\to 0$) confirms is approximately the case at small $k$.

Equation~\eqref{eq:fisher-col} is a \emph{column-wise generalized Fisher-weighted average}~\citep{matena2022merging}: task~$t$ contributes to merged column~$j$ in proportion to the squared $\ell_2$-mass of its cleaned task vector on column~$j$. Trimming depresses $\|\tilde\Delta_t[:,j]\|$ on columns where task~$t$ barely moved the weights, shifting column-$j$ routing mass to whichever task(s) did update column~$j$; in the limit of perfectly disjoint cleaned supports each $P_t$ becomes a projector onto $\{j:\|\tilde\Delta_t[:,j]\|>0\}$, sending on/off mass to infinity. The proposition also clarifies why sign-election is redundant on top of aggressive trimming: once column supports approach disjointness, little sign conflict remains to resolve.

We quantify this picture with a simple mass-fraction statistic. For each task $t$ and each layer, we define the ``active column'' set $\mathcal{S}_t$ as the top-$10\%$ of input columns by $\|\Delta_t[:,j]\|_2$ (the raw, un-trimmed task-vector's column norm). We then compute
\[
  \mathrm{on}_t(k) = \frac{\|P_t(k)[:,\mathcal{S}_t]\|_F^2}{\|P_t(k)\|_F^2},
  \qquad
  \mathrm{off}_t(k) = \operatorname*{mean}_{t'\neq t}\frac{\|P_t(k)[:,\mathcal{S}_{t'}]\|_F^2}{\|P_t(k)\|_F^2}.
\]
A ``uniform'' router ($P_t \propto I$, indifferent to task) would give $\mathrm{on}=\mathrm{off}=0.1$ and ratio $\mathrm{on}/\mathrm{off}=1.0$. ACTMat ($k{=}1$) already achieves $\mathrm{on}/\mathrm{off}=1.24$, and this ratio rises monotonically as we trim: from $1.24$ at $k{=}1$ to $1.42$ at $k{=}0.3$ (SCALE's default) and up to $2.24$ at $k{=}0.05$ (Figure~\ref{fig:routing}). The on-column mass grows from $0.099$ to $0.132$ while the off-column mass \emph{drops} from $0.080$ to $0.059$. In other words, trim-and-elect preprocessing makes each $P_t$ behave more like a task-specific projector---\emph{exactly} the inductive bias that benefits a multi-task merge. Adding sign-election on top of trim moves these numbers only slightly (on-mass $+0.005$ at $k{=}0.3$), mirroring the near-tie of trim-only and trim+sign in Table~\ref{tab:ablation}. This is our most direct empirical evidence that SCALE's advantage is a property of the routing matrix, not a property of covariance approximation quality.

\begin{figure}[t]
\centering
\includegraphics[width=\linewidth]{../results/routing_diagnostic.png}
\caption{\textbf{Task-routing diagnostic.} For each task $t$ and each of the 37 linear layers in the visual encoder, we form the routing matrix $P_t(k) = \hat C_t(k)(\sum_s \hat C_s(k) + \rho\bar\lambda I)^{-1}$ and measure what fraction of its Frobenius mass is concentrated on task-$t$'s active columns (top-$10\%$ of per-column $\|\Delta_t\|$). \textbf{(a)}~As the keep-fraction $k$ decreases, the on-target mass \emph{grows} and the off-target mass \emph{shrinks} relative to the uniform-routing null baseline of $0.1$. \textbf{(b)}~The on/off ratio rises from $1.24$ at $k{=}1$ (ACTMat) to $1.42$ at $k{=}0.3$ (SCALE default) to $2.24$ at $k{=}0.05$---a direct measure of how much more task-specific each $P_t$ becomes. This is what SCALE-Merge's ``cleaner covariances'' are doing, in the only sense that the RegMean solve cares about.}
\label{fig:routing}
\end{figure}

\subsection{Ablation study}

Table~\ref{tab:ablation} isolates the two SCALE-Merge preprocessing steps, trimming and sign-election, relative to the ACTMat baseline ($k{=}1$, no sign election). Both steps individually improve over ACTMat by 4--6 points. On our 8-task CLIP benchmark, magnitude-trimming is the \emph{dominant} contributor: trim-only at $k{=}0.3$ achieves $84.5\%$, a $6.5$ point gain. Sign-election alone (with $k{=}1$) yields $82.0\%$, a $4$ point gain. Combining both does not produce additional gain in our setting---a finding consistent with the analysis of \citet{yadav2023tiesmerging}, who observe that once the small-magnitude entries are removed, most remaining sign-conflicts are sparse and therefore already handled by the covariance reweighting inside ACTMat's closed form.

\begin{table}[t]
\centering
\caption{Ablation on the two SCALE-Merge components on the 8-task benchmark.}
\label{tab:ablation}
\footnotesize
\begin{tabular}{lccccc}
\toprule
Method & Trim ($k$) & Sign election & Ridge $\rho$ & Avg acc (\%) \\
\midrule
ACTMat (baseline)          & 1.0 & No  & $10^{-4}$ & 78.0 \\
+ Sign election only       & 1.0 & Yes & $10^{-4}$ & 82.0 \\
+ Trim only ($k{=}0.2$)    & 0.2 & No  & $10^{-4}$ & 82.8 \\
+ Trim only ($k{=}0.3$)    & 0.3 & No  & $10^{-4}$ & \textbf{84.5} \\
SCALE (trim+sign, $k{=}0.3$) & 0.3 & Yes & $10^{-4}$ & 82.8 \\
\bottomrule
\end{tabular}
\end{table}

\paragraph{Per-layer ablation: where does SCALE's gain originate?} To isolate which parts of the visual encoder benefit from trim-and-elect, we apply SCALE to a filtered subset of linear layers and ACTMat to the rest (Figure~\ref{fig:layerabl}). By \emph{layer type}, merging only the four MLP weight groups with SCALE recovers $83.4\%$ (vs.\ $84.5\%$ for SCALE-everywhere and $78.0\%$ for ACTMat), whereas an attention-only SCALE reaches just $78.9\%$---barely above ACTMat. MLP-in ($81.6\%$) and MLP-out ($80.3\%$) each individually outperform either attention component. By \emph{block depth}, SCALE on the early third of the encoder (blocks 0--3) alone recovers $83.2\%$, while the middle ($79.0\%$) and late ($77.6\%$) bands recover little. Combined: SCALE-Merge's gain over ACTMat is concentrated in \emph{early MLP layers}, consistent with prior observations that MLPs store most of a transformer's task-specific knowledge \citep{geva2021transformer,panigrahi2023task}.

\begin{figure}[t]
\centering
\includegraphics[width=\linewidth]{../results/layer_ablation.png}
\caption{\textbf{Per-layer ablation.} SCALE applied to a filtered subset of linear layers, ACTMat applied to the rest. \textbf{Left}: by layer type; MLP-only SCALE recovers nearly all the gain, attention-only does not. \textbf{Right}: by block depth; early blocks (0--3) drive most of SCALE's benefit. Red dashed line: ACTMat-everywhere baseline ($78.0\%$).}
\label{fig:layerabl}
\end{figure}

\paragraph{Sensitivity to $k$ and $\rho$.} Sweeping $k\in\{0.05,0.1,0.2,0.3,0.5,1.0\}$, the no-sign SCALE variant scores $76.7/81.1/82.8/\mathbf{84.5}/84.4/78.0$ (last $=$ ACTMat) and the sign-election variant $76.1/79.6/81.3/82.8/82.8/82.0$. Both rise sharply from $k{=}0.05$, plateau over $k\in[0.2,0.5]$, and drop at $k{=}1.0$; no-sign is consistently better for $k>0.05$. Figure~\ref{fig:ridge} shows the ridge sweep: accuracy is flat over $\rho\in[10^{-5},10^{-2}]$ and degrades only for $\rho\leq 10^{-6}$ (conditioning) or $\rho\geq 10^{-1}$ (collapse toward simple averaging); we use $\rho{=}10^{-4}$ throughout.

\paragraph{Sign-agreement diagnostic.} Consistent with Proposition~\ref{prop:rayleigh}, the majority-sign agreement rate (fraction of tasks whose trimmed task-vector entry matches the elected sign) rises with trimming from $0.65$ at $k{=}1$ to $0.73/0.81/0.92/0.95$ at $k{=}0.5/0.3/0.1/0.05$---tasks route their large updates through near-disjoint supports, leaving little sign-conflict to resolve \citep{yadav2023tiesmerging,wang2024localizing}.

\paragraph{Scaling with the number of merged tasks.} Figure~\ref{fig:taskcount} and Table~\ref{tab:taskcount_k} show average accuracy for $T\in\{2,4,6,8\}$ (first $T$ tasks alphabetically; $k$ tuned per $T$). Two trends emerge: SCALE's lead over ACTMat \emph{widens} with $T$ ($0/+2.3/+4.4/+6.5$~pp at $T{=}2/4/6/8$), and the optimal $k^\star$ \emph{decreases} from $1.0$ to $0.3$---trim earns its keep when many tasks compete for the same parameters. Since $k{=}1$ recovers ACTMat exactly, SCALE never hurts when $k$ is validated. To rule out ordering artifacts, we re-ran $T{=}4$ on three random 4-task subsets ($84.3{\pm}6.1\%$ SCALE vs.\ $81.5{\pm}6.5\%$ ACTMat; $+2.8$pp) and $T{=}6$ on $5$ random 6-task subsets from $\binom{8}{6}{=}28$ ($84.7{\pm}4.3\%$ vs.\ $76.3{\pm}5.8\%$; $+8.4$pp)---confirming the trend is not a prefix artifact.

\begin{table}[t]
\centering
\caption{Task-count ablation. SCALE with $k$ chosen per $T$ from $\{0.1, 0.2, 0.3, 0.5, 1.0\}$; the optimal $k^\star$ \emph{decreases} with $T$ because more tasks induce more interference. Other methods tuned per their standard hyperparameters.}
\label{tab:taskcount_k}
\footnotesize
\begin{tabular}{lcccc}
\toprule
Method & $T{=}2$ & $T{=}4$ & $T{=}6$ & $T{=}8$ \\
\midrule
Task Arithmetic             & 88.4 & 75.5 & 73.2 & 65.5 \\
TIES-Merging                & 92.9 & 76.2 & 71.2 & 64.0 \\
ACTMat                      & 81.2 & 82.2 & 84.4 & 78.0 \\
SCALE-Merge (best $k^\star$)  & 81.2 & 84.5 & 88.8 & 84.5 \\
\hspace*{1em}SCALE $k^\star$ & 1.0 & 0.2 & 0.3 & 0.3 \\
\bottomrule
\end{tabular}
\end{table}

\begin{figure}[t]
\centering
\begin{minipage}{0.48\linewidth}
\includegraphics[width=\linewidth]{../results/ridge_sweep.png}
\caption{Average test accuracy of SCALE-Merge as a function of the ridge $\rho$. Accuracy is flat over four orders of magnitude, indicating robustness.}
\label{fig:ridge}
\end{minipage}\hfill
\begin{minipage}{0.48\linewidth}
\includegraphics[width=\linewidth]{../results/taskcount.png}
\caption{Accuracy vs.\ number of merged tasks $T$. SCALE-Merge's lead over ACTMat grows with $T$.}
\label{fig:taskcount}
\end{minipage}
\end{figure}

\paragraph{Seed robustness.} Repeating the full 8-task fine-tuning under three independent seeds ($\{42,123,456\}$) and re-merging with unchanged hyperparameters, ACTMat averages $78.5{\pm}0.4\%$, SCALE (trim-only, $k{=}0.3$) $84.6{\pm}0.2\%$ ($+6.05{\pm}0.39$~pp), and SCALE (trim+sign) $82.7{\pm}0.4\%$ ($+4.18{\pm}0.41$~pp). Trim-only $>$ trim+sign $>$ ACTMat on every seed, so the Table~\ref{tab:ablation} ordering is seed-robust (see Appendix~\ref{app:seed}).

\subsection{Cross-backbone validation: CLIP ViT-L/14}
\label{sec:vitl14}

We repeat the core comparison on the $\sim$3$\times$-larger CLIP ViT-L/14 backbone ($d_\text{in}\le 4096$, same fine-tuning recipe; Table~\ref{tab:vitl14}). SCALE remains best ($92.1\%$ at $k^\star{=}0.5$, $+0.9$pp over ACTMat); the smaller gap vs.\ ViT-B/32 reflects that ACTMat already closes most of the ZS$\to$UB gap on this backbone. The optimal $k$ rises from $0.3$ to $0.5$---larger $d_\text{in}$ tolerates less aggressive trimming. SCALE wins on $5/8$ tasks (largest gains CIFAR-100 $+2.9$, EuroSAT $+2.3$, DTD $+1.8$); the three losses are all within $0.7$pp. All four sparsification-family baselines (TA $79.9$, TIES $78.9$, DELLA $79.3$, Breadcrumbs $80.1$) cluster in a $1.2$-pp band and trail ACTMat by $\geq 11$pp, mirroring the ViT-B/32 ordering.

\begin{table}[t]
\centering
\caption{\textbf{Cross-backbone validation on CLIP ViT-L/14.} Top-1 accuracy (\%) per task, and average across the 8 tasks, merging 8 fine-tuned visual encoders. Best merged avg in bold.}
\label{tab:vitl14}
\footnotesize
\setlength{\tabcolsep}{4pt}
\resizebox{\linewidth}{!}{
\begin{tabular}{lccccccccc}
\toprule
Method & MNIST & CIFAR10 & CIFAR100 & SVHN & FashionM & EuroSAT & GTSRB & DTD & \textbf{Avg} \\
\midrule
Pretrained (zero-shot)  & 55.2 & 95.3 & 73.3 & 53.7 & 65.7 & 52.0 & 48.8 & 52.4 & 62.1 \\
Individual (UB)         & 99.8 & 99.1 & 93.6 & 98.1 & 95.8 & 99.0 & 99.4 & 81.7 & 95.8 \\
\midrule
Simple Averaging        & 89.2 & 98.0 & 85.0 & 76.4 & 84.6 & 68.4 & 59.4 & 57.7 & 77.3 \\
Task Arithmetic ($\alpha{=}0.2$)    & 96.3 & 98.1 & 84.5 & 82.0 & 87.6 & 69.5 & 63.3 & 57.8 & 79.9 \\
TIES-Merging ($k{=}0.2,\alpha{=}0.5$) & 94.8 & 98.5 & 86.6 & 81.9 & 87.4 & 66.9 & 59.7 & 55.4 & 78.9 \\
DELLA ($p_\ell{=}0.1,p_h{=}0.8,\alpha{=}0.5$) & 96.4 & 98.3 & 86.3 & 84.2 & 88.0 & 65.7 & 60.9 & 54.8 & 79.3 \\
Breadcrumbs ($s{=}0.85,l{=}0.005,\alpha{=}0.5$) & 97.5 & 98.0 & 83.9 & 83.7 & 88.6 & 67.0 & 64.3 & 57.4 & 80.1 \\
ACTMat                  & \textbf{99.6} & 98.3 & 86.2 & \textbf{93.9} & 93.5 & 94.8 & \textbf{94.7} & 68.4 & 91.2 \\
\textbf{SCALE-Merge (ours, $k^\star{=}0.5$)} & 99.5 & \textbf{98.9} & \textbf{89.1} & 93.3 & \textbf{94.4} & \textbf{97.1} & 94.1 & \textbf{70.2} & \textbf{92.1} \\
\bottomrule
\end{tabular}}
\end{table}

\section{Discussion}

\paragraph{Connections to MaTS and a matrix-free CG variant.} Our method inherits the linear-system structure from MaTS \citep{tam2024mats}: SCALE-Merge picks the task parameter subspace $Q_t$ to be the column-space of the cleaned task vector $\tilde\Delta_t$, which is a ``data-free, sign-aware, sparsified'' choice of covariance in MaTS terminology. Following \citet{tam2024mats}, we implement a matrix-free conjugate-gradient solver that never materialises the $d_\text{in}{\times}d_\text{in}$ gram matrix $\sum_t \hat C_t$---each CG matvec applies $v\mapsto \sum_t\tilde\Delta_t^\top(\tilde\Delta_t v)+\rho\bar\lambda\,v$ directly. With a $10^{-6}$ residual cap of $50$ iterations, CG recovers the direct-solve accuracy to within $0.7$~pp on ViT-B/32 ($83.8$ vs.\ $84.5$) and within $0.2$~pp on ViT-L/14 ($91.8$ vs.\ $92.1$), with mean $39$--$47$ CG iterations per layer. At the scales we study, the direct $\mathrm{LU}$ solve is faster (ViT-B/32: $0.56$s vs.\ $6.3$s; ViT-L/14: $1.4$s vs.\ $28.5$s) and the memory savings are negligible because the $d_\text{in}{\times}d_\text{in}$ covariance ($\leq 64$\,MB at $d_\text{in}{=}4096$) is dominated by the $8$-task checkpoint mass on GPU ($6.7$/$22.4$\,GiB peak); CG becomes useful once $d_\text{in}$ exceeds what the dense solve fits in memory.

\paragraph{Limitations and cost.} SCALE-Merge inherits from ACTMat layer-wise independence and reliance on the gradient-descent regime under which $\Delta_t^\top\Delta_t$ approximates $C_t$. The trim step is most effective when top-$k$ supports are near-disjoint; when tasks compete for the \emph{same} parameters, sign-election becomes the dominant factor. Computationally, SCALE-Merge merges the 8-task ViT-B/32 benchmark in $0.45$s on one H100 (vs.\ $0.44$s ACTMat, $0.24$s TIES, $0.07$s Task Arithmetic) and is orders of magnitude cheaper than data-based RegMean variants.

\section{Conclusion}

We introduced SCALE-Merge, a data-free method that unifies interference-resolving preprocessing (TIES) with activation-matching merging (ACTMat/RegMean). It obtains a meaningful improvement over every data-free and data-based baseline we test, including a RegMean oracle given true per-task activation covariances. Our diagnostic experiments show that this gain does \emph{not} come from a better approximation of the activation covariance---trimming in fact degrades that approximation---but from a better \emph{effective reweighting} of the merge linear system, which our task-routing-matrix diagnostic measures directly: SCALE nearly doubles ACTMat's on-target routing concentration. We view this as an instance of a broader principle: \emph{the covariance in an activation-matching merge is less a statistical quantity to be estimated accurately than an inductive bias to be designed.} Further gains should be possible by combining SCALE-Merge with adaptive covariance estimators \citep{aceemrging2026} and low-rank factorizations of $\Delta_t$.

% References
\bibliographystyle{iclr2026_conference}
\bibliography{refs}

\newpage
\appendix

\section{Proof of Proposition~\ref{prop:rayleigh}}
\label{app:proof}

\begin{proof}
From the RegMean/MaTS closed form $\ws = \bigl(\sum_t W_t \hat C_t\bigr) M^{-1}$, extracting column $j$ gives
$\ws[:, j] = \sum_t W_t\,(\hat C_t M^{-1} e_j) = \sum_t W_t\, P_t[:, j]$,
which proves the first claim with $P_t[:, j] = \hat C_t M^{-1} e_j$.

For the diagonal reduction, assume each $\hat C_t$ is diagonal. Then its diagonal entries are $\hat C_t[j, j] = \|\tilde\Delta_t[:, j]\|_2^2$, and $M = \sum_s \hat C_s + \rho\bar\lambda I$ is also diagonal with entries $M[j, j] = \sum_s \|\tilde\Delta_s[:, j]\|_2^2 + \rho\bar\lambda$. Therefore $M^{-1}e_j = e_j / M[j, j]$ and
\[
P_t[:, j]
\;=\; \hat C_t\,\frac{e_j}{M[j, j]}
\;=\; \frac{\|\tilde\Delta_t[:, j]\|_2^2}{\sum_s \|\tilde\Delta_s[:, j]\|_2^2 + \rho\bar\lambda}\, e_j,
\]
whose $j$-th coordinate is the right-hand side of~\eqref{eq:fisher-col}.

A sufficient condition for each $\hat C_t$ to be diagonal is that, for every pair of distinct input columns $i\neq j$, the columns $\tilde\Delta_t[:, i]$ and $\tilde\Delta_t[:, j]$ have disjoint row-supports (no output row is nonzero in both). The TIES trim keeps only a $k$-fraction of the $d_\text{in}\cdot d_\text{out}$ entries of $\Delta_t$ by magnitude; under a simple random-support surrogate the expected overlap of two columns' supports is a $k^2$-fraction of rows, which is $\leq 0.09$ at our default $k=0.3$ and vanishes as $k\to 0$. When the assumption holds only approximately, $P_t[:, j]$ picks up an $O(\|\hat C_t - \mathrm{diag}(\hat C_t)\|_{\mathrm{op}} / M[j, j])$ off-diagonal leakage from neighboring columns; \S\ref{sec:analysis} shows this leakage is the dominant failure mode at $k=1$ (ACTMat), where the on/off-target routing ratio collapses to $1.24$.
\end{proof}

\paragraph{Connection to Fisher merging.} Equation~\eqref{eq:fisher-col} reproduces the form of Fisher-weighted averaging~\citep{matena2022merging} when the task-$t$ Fisher diagonal at column $j$ is taken to be $\|\tilde\Delta_t[:, j]\|_2^2$---the squared-gradient data-free approximation of~\citet{pascanu2014revisiting} applied column-wise to the cleaned task vector. SCALE-Merge strictly generalizes this diagonal Fisher reduction: the full $\hat C_t$ carries off-diagonal cross-column couplings that a purely column-wise Fisher average discards, and it is precisely these couplings that drive the routing-matrix concentration we measure in Figure~\ref{fig:routing}.

\section{Seed-robustness details}
\label{app:seed}

We repeat the 8-task fine-tuning recipe of~\S\ref{sec:exp} from scratch under each of three seeds $\{42, 123, 456\}$, using the same hyperparameters. Individual-task accuracies vary by at most $\pm 0.6$~pp across seeds (largest variation: DTD at $74.4$--$75.0\%$). We then apply ACTMat, SCALE (trim-only, $k{=}0.3$, $\rho{=}10^{-4}$), and SCALE (trim+sign, same $k$, $\rho$) to each seed's checkpoints; Table~\ref{tab:seed_appendix} lists per-seed averages.

\begin{table}[h]
\centering
\caption{Per-seed 8-task average accuracy (\%). See \texttt{results/seed\_robustness.json} for per-task numbers.}
\label{tab:seed_appendix}
\footnotesize
\begin{tabular}{lcccc}
\toprule
Method & seed 42 & seed 123 & seed 456 & Mean$\pm$std \\
\midrule
ACTMat              & 78.43 & 78.06 & 79.07 & $78.5{\pm}0.4$ \\
SCALE (trim-only)   & 84.26 & 84.66 & 84.80 & $84.6{\pm}0.2$ \\
SCALE (trim+sign)   & 82.20 & 82.80 & 83.10 & $82.7{\pm}0.4$ \\
\bottomrule
\end{tabular}
\end{table}

\end{document}
